% Baccalauréat malgache, série D, 2012 : énoncé.
\TCDBacAttribution{D}{2012}
\noindent \underline{\textsc{Exercice I}} (5 points) 
      \begin{enumerate}
            \item Calculer $(2 - i)^2$. 
                  En déduire la résolution dans $\mathbb{C}$ l'équation $(E) : iz^2 - iz + 1 + i = 0$.
            \item Soit $P$ le polynôme de variable complexe $z$ définie par : $P(z) = z^3 - (4+i)z^2 + (7+i)z - 4$.
                  \begin{enumerate}
                        \item Montrer que l’équation $P(z) = 0$ admet une solution 
                              réelle positive $\alpha$ que l’on déterminera.
                        \item Déterminer les nombres $a$ et $b$ tels que : $P(z) = (z - 1) (z - 2 - 2i) (az + b)$.
                  \end{enumerate}
            \item Dans le plan complexe muni d’un repère orthogonal direct $R(O;\vec{u},\vec{v})$ on consid\`ere les
                  trois points $A$, $B$ et $C$, d’affixes respectives : $1$ ; $2 + 2i$ et $1 - i$.
                  \begin{enumerate}
                        \item Donner la forme trigonométrique de $U = \frac{2+2i}{1-i}$.
                        \item En déduire la nature du triangle $OBC$.
                  \end{enumerate}
            \item Soit $S$ la similitude plane directe telle que :
                  $\displaystyle{ \begin{multiform} x' = x + y - 1 \\ y' = -x + y + 2 \;. \end{multiform} }$
                  Déterminer l’expression complexe de $S$ et préciser ses éléments caractéristiques.
      \end{enumerate}
%%%%%%%%%%%%%%%%%%%%%%%%%%%%%%%%%%%%%%%%%%%%%%%%%%%%%%%%%%%%%%%%%%%%%%%%%%%%%%%%%%%%%%%%%%%%%%%
%%%%%%%%%%%%%%%%%%%                   %%%%%%%%%%%%%%%%%%%                   %%%%%%%%%%%%%%%%%%%
%%%%%%%%%%%%%%%%%%%%%%%%%%%%%%%%%%%%%%%%%%%%%%%%%%%%%%%%%%%%%%%%%%%%%%%%%%%%%%%%%%%%%%%%%%%%%%%
\noindent \underline{\textsc{Exercice II}} (5 points) \newline
      Un bassin contient 10 poissons, dont 2 carpes, 3 tanches et 5 gardons.
      On prélève au hasard et simultanément 3 poissons du bassin à l'aide d'un filet. Chaque poisson a la même
      probabilité d’être pris par le filet.
      \begin{enumerate}
            \item Calculer la probabilité de chacun des évènements suivants :
                  \begin{description}
                        \item[$E_1$ :] « ne pêcher aucune tanche »,
                        \item[$E_2$ :] « pêcher au moins une carpe »,
                        \item[$E_3$ :] « pêcher exactement 2 gardons ».
                  \end{description}
                  On répète quatre fois de suite cette épreuve d’une manière indépendante.
            \item Soit $X$ la variable aléatoire égale au nombre de réalisations de l’évènement $E_3$.
                  \begin{enumerate}
                        \item Donner la loi de probabilité de $X$.
                        \item Calculer l’espérance mathématique $E(x)$ et la variance $V(x)$.
                        \item Calculer la probabilité $P(x \leq 3)$.
                  \end{enumerate}
            (N.B : on donnera les résultats sous forme de fraction irréductible).
      \end{enumerate}
%%%%%%%%%%%%%%%%%%%%%%%%%%%%%%%%%%%%%%%%%%%%%%%%%%%%%%%%%%%%%%%%%%%%%%%%%%%%%%%%%%%%%%%%%%%%%%%
%%%%%%%%%%%%%%%%%%%                   %%%%%%%%%%%%%%%%%%%                   %%%%%%%%%%%%%%%%%%%
%%%%%%%%%%%%%%%%%%%%%%%%%%%%%%%%%%%%%%%%%%%%%%%%%%%%%%%%%%%%%%%%%%%%%%%%%%%%%%%%%%%%%%%%%%%%%%%
\noindent \underline{\textsc{Probleme}} 
      \begin{description}
            \item[I.]
                  Soit $g$ la fonction numérique définie sur $]0 ; +\infty [$ par : $g(x) = x^2 +2 - 2 \ln x$
                  où $\ln$ désigne la fonction logarithme népérien.
                  \begin{enumerate}
                        \item Calculer $g'(x)$ puis étudier le sens de variation de $g$ sur $]0 ; +\infty[$. 
                              (on ne demande pas les limites).
                        \item Calculer $g(1)$. 
                              En déduire que $g$ est strictement positif sur $]0 ; +\infty[$.
                  \end{enumerate}
            \item[II.]
                  Soit $f$ la fonction numérique définie sur l’intervalle $] 0 ; +\infty[$ par :
                  $f(x) = x - 1 + \frac{2 \ln x}{x}$.
                  On désigne par $(\mathscr{C})$ sa courbe représentative dans un repère orthonormé
                  d'unité graphique 2~cm.
                  \begin{enumerate}
                        \item \begin{enumerate}
                                    \item Calculer $\displaystyle{\lim_{x\rightarrow 0^+}f(x)}$. 
                                          Interpréter graphiquement le résultat.
                                    \item Calculer $\displaystyle{\lim_{x\rightarrow +\infty}f(x)}$.
                              \end{enumerate}
                        \item \begin{enumerate}
                                    \item Montrer que la droite $(\Delta)$ d'\'equation $y=x-1$ est une asymptote oblique 
                                          \`a la courbe $(\mathscr{C})$ au voisinage de $+\infty$.
                                    \item Etudier la position de $(\mathscr{C})$ par rapport \`a $(\Delta)$.
                              \end{enumerate}
                        \item \begin{enumerate}
                                    \item Montrer que pour tout $x\in ]0;+\infty[$, $f'(x) = \frac{g(x)}{x^2}$
                                          o\`u $f$ d\'esigne la fonction d\'eriv\'ee de $f$.
                                    \item Justifier que $f'(x)$ \`a le m\^eme signe que $g(x)$ suivant les valeurs de $x$.
                                    \item Dresser le tableau de variation de $f$ sur $]0;+\infty[$.
                              \end{enumerate}
                        \item \begin{enumerate}
                                    \item Ecrire une \'equation de la tangente $(T)$ \`a $(\mathscr{C})$ 
                                          au point d'abscisse $x_0 = 1$.
                                    \item Tracer $(T)$, $(\Delta)$ et $(\mathscr{C})$ dans le m\^eme rep\`ere.
                              \end{enumerate}
                        \item Soit $h$ la fonction d\'efinie sur $]0;+\infty[$ par : $h(x) = \frac{1}{2} (\ln x)^2$.
                              \begin{enumerate}
                                    \item D\'eterminer $h'(x)$. En d\'eduire l'expression 
                                          $I(\alpha) = \int_{1}^{\alpha} f(x) \dx$ o\`u $\alpha > 1$, en fonction de $\alpha$.
                                    \item Calculer, en $\mathrm{cm}^2$, l'aire $A$ du domaine plan limit\'e par
                                          la courbe $(\mathscr{C})$, l'axe des abscisses et les droites 
                                          d'\'equations respectives $x=1$ et $x=e$.
                              \end{enumerate}
                        On donne $e \simeq 2,7$.
                  \end{enumerate}
      \end{description}












